% 本文件由 LaTeX 排版 Agent 根据有效 translation.json 自动生成。
% author=Apocrypha; work_id=0826; title=十二宗徒之一达陡行传

\chapter{\WorkTitle{十二宗徒之一达陡行传}}

% structure: p0001 source=h1 target=omitted reason=page-title-already-rendered-by-parent

\Person{肋拜乌斯}（即\Person{达陡}）是\Place{厄得撒}城人——该城是\Place{奥斯若恩}的首府，位于\Place{亚美尼亚叙利亚}内陆。他生为希伯来人，精通圣经，极有学识。在\Person{若翰}洗者时期，他到\Place{耶路撒冷}朝拜；听了他的宣讲，目睹了他天使般的生活，便受了洗，他的名被称为\Person{达陡}。他见到基督的显现、祂的教导和奇妙作为，便跟随了祂，成了祂的门徒；祂拣选他作为十二宗徒之一，按福音作者\Person{玛窦}和\Person{马尔谷}的记载，他是第十位宗徒。\par

那时，\Place{厄得撒}城的总督名叫\Person{阿布加尔}。基督的名声、祂所行的奇事及祂的教导传遍各地，\Person{阿布加尔}听闻后大为惊异，渴望见到基督，却不能离开他的城和政府。在受难前后，犹太人设谋迫害的那段日子，\Person{阿布加尔}被不治之症缠身，便通过信使\Person{阿纳尼亚斯}给基督送了一封信，内容如下：\Quote{致被称为基督的\Person{耶稣}，\Place{厄得撒}地区的总督\Person{阿布加尔}，不配的奴隶。我所听到的你所行的众多奇事，就是你治愈盲人、瘸子、瘫痪者，并治好所有附魔的人。因此我恳求你的良善，甚至到我们这里来，躲避那些邪恶的犹太人因嫉妒而对你发起的阴谋。我的城虽小，却足够容纳我们两人。}\Person{阿布加尔}吩咐\Person{阿纳尼亚斯}要仔细记下基督的相貌、身高、头发，总而言之，一切细节。\par

\Person{阿纳尼亚斯}去了，递交了信，并仔细端详基督，却无法在心中确定祂的样貌。基督知道人心，便请求沐浴；有人递给祂一条毛巾。祂洗完澡，用毛巾擦了脸。祂的肖像便印在了毛巾上，祂将毛巾交给\Person{阿纳尼亚斯}，说：\Quote{拿去这个，并给那打发你来的人带回这话：愿你和你城平安！因为我正是为此而来，为世界受苦、复活，并唤醒先祖。当我被接到天上之后，我会派我的门徒\Person{达陡}到你那里，他要光照你，引导你进入一切真理，你和你城都一样。}\par

\Person{阿纳尼亚斯}回来后，\Person{阿布加尔}俯伏敬拜了那肖像，在\Person{达陡}来到之前，他的病便已痊愈。\par

在基督受难、复活、升天之后，\Person{达陡}到了\Person{阿布加尔}那里；见他已康复，便向他讲述了基督降生成人的事迹，为他及他全家施了洗。他还教导了大批群众，包括希伯来人、希腊人、叙利亚人和亚美尼亚人，并因父、及子、及圣神之名给他们施洗，用圣油傅抹了他们；又让他们分享了主耶稣基督圣体圣血的无玷奥迹，并吩咐他们遵守梅瑟法律，留意宗徒们在耶路撒冷所传的话。因为他们每年都聚在一起过逾越节，他又把圣神赐给了他们。\par

\Person{达陡}与\Person{阿布加尔}一起摧毁了偶像庙宇，建立了教堂；祝圣了他的一位门徒为主教，并祝圣了长老和执事，为他们制定了圣咏和圣礼仪的规则。然后他离开了他们，前往\Place{阿弥斯}城，那是\Place{默塞加尔德人}和\Place{叙利亚人}的大都会，即\Place{美索不达米亚-叙利亚}，位于\Place{底格里斯河}畔。在安息日，他与门徒们进入犹太人的会堂；诵读法律后，大司祭对\Person{达陡}和他的门徒们说：\Quote{诸位，你们从何处来？为何在此？}\par

\Person{达陡}说：\Quote{你们无疑已听说过在\Place{耶路撒冷}关于\Person{耶稣基督}所发生的事，我们是祂的门徒，是祂所行所教的奇事的见证人，以及司祭长们如何因仇恨将祂交给\Place{犹太}总督\Person{比拉多}。\Person{比拉多}审问了祂，找不到罪状，本想释放祂；但他们却喊叫说：\InnerQuote{你若释放这人，就不是凯撒的朋友，因为他自称是君王。}他害怕了，在众人面前洗了手，说：\InnerQuote{对这人的血，我是无罪的，你们自己负责。}司祭长们回答说：\InnerQuote{祂的血归到我们和我们的子孙身上。}于是\Person{比拉多}把祂交给了他们。他们拿住祂，与兵士一起向祂吐唾沫，大大戏弄祂，然后钉了祂十字架，将祂安放在坟墓里，又严密看守，还派了兵士把守。第三天黎明前，祂复活了，将殓布留在坟墓里。祂首先显现给祂的母亲和其他妇女，而后显现给\Person{伯多禄}和\Person{若望}——他们是我同门中最早的，以后又显现给我们十二人。祂复活后多日，我们与祂一同吃喝。祂派遣我们以祂的名向万民宣讲悔改和罪过的赦免，使那些受洗并听到天国福音的人，在今世终结时成为不朽的复活；祂赐给我们权柄驱逐魔鬼，治愈一切疾病和灾殃，并使死人复活。}\par

众人听了这话，便将他们的病人和附魔的人带来。\Person{达陡}与门徒们出去，将手按在每个人身上，呼求基督的名，把他们都治好了。附魔的人在\Person{达陡}走近之前就已痊愈，魔鬼从他们身上出去了。多日以来，人们从各处跑来，观看\Person{达陡}所行的事。听了他的教导，许多人信了，并宣认自己的罪过而受了洗。\par

因此，与他们相处五年后，他建造了一座圣堂；并任命他的一位门徒为主教，又任命了司铎和执事，为他们祈祷后，便离开，巡行\Place{叙利亚}各城，教导并治愈所有病人；因此，他藉着教导，使许多城邦和国家归向了基督。于是，他同门徒一起教导、传福音，并医治病人，来到了海边的\Place{腓尼基}城\Place{贝里图斯}；在那里教导并光照了许多人之后，于八月二十一日安息了。门徒们聚集起来，以极大的尊荣埋葬了他；许多病人得到痊愈，他们便光荣圣父、圣子及圣神，至于永远。阿们。\par

\SourceRecord{Translated by Alexander Walker.；Ante-Nicene Fathers；Vol. 8.；Edited by Alexander Roberts, James Donaldson, and A. Cleveland Coxe.；Buffalo, NY: Christian Literature Publishing Co.,；1886.；<http://www.newadvent.org/fathers/0826.htm>.}
